%============================================================
% TEMPLATE — PATTERN THEORY
% Copy to chapters/chNN_xxx/pK_xxx/theory.tex
%============================================================
\subsection{Pattern Theory}
\label{ssec:chN-pK-theory}

\begin{patternbox}{Pattern K at a Glance}
% one picture + the core algorithm (5-8 lines of pseudocode) + one sentence
\end{patternbox}

\subsubsection*{From the previous pattern to this one}
% what single ingredient changes (decision, dimension, modes, meaning)

\subsubsection*{When to recognise it}
% "use when" list and "do NOT use when" list (point to the right pattern)

\subsubsection*{The framework}
% table: state / transition / base / answer / order / memory
% the invariant of the algorithm and a proof sketch (invariantbox)

\subsubsection*{Variations map}
% diagram: root node + dashed group of variations (see Chapter 1 patterns)
% bullet list: twist -> how to handle it

\subsubsection*{Code templates}
% compilable [style=cpp] with a main() that prints checked values

\subsubsection*{Common mistakes}
% warningbox per mistake, wrong vs right snippet

\begin{checklistbox}[``I'm stuck'' checklist for Pattern K]
\begin{itemize}[leftmargin=*]
    \item[$\square$] \ldots
\end{itemize}
\end{checklistbox}
